\section{Genuine generic power frames and true proper-model norm coherence}

\begin{lemma}\label{lem:fc-generic-tensor-frame-unit-norm}
\lean{EnvelopeProperty.Berkovich.SchemeLinePowerFrameNorms.genericTensorFrame,EnvelopeProperty.Berkovich.SchemeLinePowerFrameNorms.genericTensorFrame_unitSection,EnvelopeProperty.Berkovich.SchemeLinePowerFrameNorms.modelLineNorm_genericTensorFrame_unitSection}
\leanok
\uses{lem:fc-actual-tensor-frame-unit-coefficient,thm:actual-model-line-tensor-norm,lem:bundle-tensor-power-pullback}
For actual entire frames of the same original residue fibres of
$M$ and $N$ at a native point $y$, the actual native pullback
tensorator and explicit tensor comparison give an entire frame
of the original explicit tensor fibre. Its inverse unit section
is the actual generic tensor of the two original inverse unit
sections. If the original structure map is proper over a valuation
bounded by one and both original modules are invertible, the true
model norm of this inverse unit section is the product of the
two true original frame-unit model norms.
\end{lemma}
\begin{proof}
Apply the actual inverse pullback tensorator to identify the
original tensor fibre with the tensor of the same original fibres.
Follow it by the actual tensor frame, preceding these maps by
the original pulled explicit tensor comparison. The entire
inverse-unit precomposition identity and actual tensor-frame
unit-section formula identify its inverse unit section with the
actual generic pure tensor through those same native comparisons.
The genuine model-norm tensor theorem computes the norm of this
actual generic tensor as the product of the true original norms.
No coefficient, frame-preservation or norm comparison is assumed.
\end{proof}

\begin{theorem}\label{thm:fc-generic-power-frame-true-norm}
\lean{EnvelopeProperty.Berkovich.SchemeLinePowerFrameNorms.genericPowerFrame,EnvelopeProperty.Berkovich.SchemeLinePowerFrameNorms.genericPowerFrame_unitSection_succ,EnvelopeProperty.Berkovich.SchemeLinePowerFrameNorms.transitionUnit_genericPowerFrame,EnvelopeProperty.Berkovich.SchemeLinePowerFrameNorms.transitionUnit_genericPowerFrame_value,EnvelopeProperty.Berkovich.SchemeLinePowerFrameNorms.modelLineNorm_genericPowerFrame_unitSection}
\leanok
\uses{lem:fc-generic-tensor-frame-unit-norm,thm:fc-actual-power-frame-transition-units,lem:scheme-line-actual-model-unit-norm}
Any actual entire original residue frame $e$ of $M$ induces actual
entire frames $e_n$ on the same original residue fibres of every
explicit $M^{\otimes n}$. The zeroth frame is the actual native
structure-unit pullback comparison; successors use the genuine
generic tensor-frame construction. For two original frames,
their genuinely derived transition units satisfy $u_n=u_1^n$,
so every native residue valuation gives $c_n=c_1^n$. Under the
proper bounded-base hypotheses and original invertibility,
\[
 \|e_n^{-1}(1)\|_{M^{\otimes n},y}
   =\|e^{-1}(1)\|_{M,y}^{\,n}.
\]
All powers, including zero, use actual original fibres and true
model norms; no power compatibility identity is an input.
\end{theorem}
\begin{proof}
Construct the frames by the stated recursion on actual original
module isomorphisms. The preceding inverse-unit formula gives
the actual generic tensor at every successor. To compare two
frames, common-precomposition invariance removes the same actual
pulled tensor comparisons; the actual tensor-frame transition
identity multiplies the preceding transition by the degree-one
transition. Induction, beginning with identical degree-zero
structure frames, proves the actual unit equality. Applying the
valuation proves its exact value equality. For the true norm,
the original structure-unit norm formula gives value one on the
zeroth frame's inverse unit section. At a successor the preceding
lemma gives multiplication by the degree-one frame-unit norm.
The induction hypothesis and the power recurrence prove the
asserted genuine norm equality.
\end{proof}

\begin{theorem}\label{thm:fc-proper-center-power-frame-coherence}
\lean{EnvelopeProperty.Berkovich.SchemeLinePowerFrameNorms.centerPowerFrame,EnvelopeProperty.Berkovich.SchemeLinePowerFrameNorms.centerFrameGenericIso_centerPowerFrame,EnvelopeProperty.Berkovich.SchemeLinePowerFrameNorms.modelLineNorm_centerPowerFrame_unitSection}
\leanok
\uses{thm:fc-generic-power-frame-true-norm,lem:actual-center-generic-tensor-frame,lem:scheme-line-actual-generic-frame-naturality,lem:scheme-line-actual-model-unit-norm}
Every actual original valuative-center frame of $M$ induces actual
entire center frames of all the same original explicit powers.
Their entire generic frames equal exactly the preceding power
frames formed from the original degree-one generic frame. Their
true proper-model inverse-unit norms are all one. These are
whole original frame identities derived from actual native tensor
and pullback comparisons; no preservation or power-coherence
certificate is supplied.
\end{theorem}
\begin{proof}
At zero choose the true native unit pullback comparison on the
actual proper valuation ring. At a successor use the genuine
center tensor frame and the original pulled explicit tensor
comparison. For the whole generic identity, induct on the power.
At zero, genuine structure-unit compatibility with native pullback
composition and the actual valuative triangle identifies its
entire generic frame with the native residue unit comparison.
At a successor, generic-frame naturality moves the actual
explicit tensor comparison past the proper-center generic
comparison. The genuine center tensor-frame theorem identifies
the remaining entire generic tensor frame, and the induction
hypothesis identifies its preceding factor with the actual
previous generic power frame. This proves the whole equality.
Finally any actual center frame computes the true model norm.
Its entire generic inverse unit section has original coefficient
one by the frame inverse identity, so its actual native residue
valuation, and hence its true model norm, is one.
\end{proof}
